\pnum{5.24}{Wavefunctions in a periodic delta-function comb}
Inside a cell $0<x<a$ the general solution
is $\psi(x)=A\sin(kx)+B\cos(kx)$.
Continuity and the Bloch condition imply
$A\sin(ka)=[\ee^{\ii qa}-\cos(ka)]B$.
For $\sin(ka)\ne0$, substitute and
use $\sin(ka)\cos(kx)-\cos(ka)\sin(kx)
=\sin[k(a-x)]$:
\begin{equation}
\boxed{\psi(x)=C\{
\sin(kx)+\ee^{-\ii qa}\sin[k(a-x)]\}},
\qquad 0\le x\le a,
\end{equation}
where $C=A/[\ldots]$ absorbs an
unimportant constant and phase.
Equivalently this follows by choosing
$C=B\ee^{\ii qa}/\sin(ka)$.
At a band boundary $z=ka=j\pi$,
the ratio used in this derivation is
singular. The correct state has zeros
at every delta spike:
\begin{equation}
\boxed{\psi(x)\propto
\sin(j\pi x/a)}
\end{equation}
extended globally as a standing wave.
Its derivative is continuous across
each spike because the jump
$2m\alpha\psi/\sh^2$ vanishes there.
The Bloch phase at this boundary is
$\ee^{\ii qa}=(-1)^j$.

\pnum{5.25}{Lower edge of the first repulsive delta-comb band}
Define $z=ka$ and
$\beta=m\alpha a/\sh^2$.
The Bloch dispersion equation is
\begin{equation}
\boxed{\cos(qa)=f(z)=
\cos z+\beta\frac{\sin z}{z}}.
\end{equation}
For $\beta=10$, $f(0)=11$,
which is forbidden because $|\cos qa|\le1$.
The lowest band starts at the first
positive root of
$f(z)=1$, which can be found with
bisection or Brent's method.
Numerically, $\boxed{z_0=2.627675433\ldots}$.
The energy is $E=\sh^2z^2/(2ma^2)$.
Since $\alpha/a=1$ eV and
$\beta=m\alpha a/\sh^2=10$,
$\sh^2/(2ma^2)=\alpha/(2\beta a)=0.05$ eV.
Therefore
\begin{equation}
\boxed{E_{\rm band,bottom}=
0.05z_0^2\ {\rm eV}\simeq0.345\ {\rm eV}}
\end{equation}
(three significant digits).

\pnum{5.26}{Attractive delta comb and negative-energy Bloch bands}
For an attractive lattice set
$\beta=-|\beta|$ in the same
positive-energy dispersion function:
\begin{equation}
f(z)=\cos z-|\beta|\frac{\sin z}{z},
\qquad z=ka>0.
\end{equation}
For negative energies, let
$\kappa=\sqrt{-2mE}/\sh>0$
and $t=\kappa a$. Analytic continuation
$k=\ii\kappa$ gives
\begin{equation}
\boxed{\cos(qa)=
\cosh t-|\beta|\frac{\sinh t}{t}}.
\end{equation}
Allowed states have the right-hand
side in $[-1,1]$. For $|\beta|=1.5$,
$f(E=0)=1-1.5=-0.5$, so the
lowest band crosses $E=0$, containing
both negative- and positive-energy
states. As $E\to-\infty$,
$\cosh t-|\beta|\sinh t/t\to\infty$
and negative energies become forbidden.
There is therefore a finite lower edge
for the bound portion of the band.
A periodic chain of $N$ cells has
$\boxed{N\text{ spatial Bloch states
per allowed band}}$ with periodic
boundary conditions. Only some may
have $E<0$ in this parameter regime;
the rest lie in its positive-energy
continuation.

\pnum{5.27}{Bloch momentum reversal and energy degeneracies}
For a ring of $N$ identical cells,
$q_n=2\pi n/(Na)$, with
$n=0,\ldots,N-1$.
The dispersion relation depends
on $q$ through $\cos(qa)$.
Hence
\begin{equation}
E(q)=E(-q)
\end{equation}
within any fixed band.
The modes $n$ and $N-n$ have the
same cosine and normally form a pair
of degenerate but independent traveling
Bloch states. Exceptions occur when
the two modes coincide modulo the
reciprocal lattice vector:
\begin{equation}
q\equiv-q\pmod{2\pi/a}.
\end{equation}
This means $2n\equiv0\pmod N$.
Thus
\begin{equation}
\boxed{n=0\text{ always};\qquad
n=N/2\text{ also if }N\text{ is even}}.
\end{equation}
Those time-reversal invariant momenta
are not forced into distinct $\pm q$
pairs. At zero and zone-edge
momenta there can be nondegenerate
band-edge eigenfunctions.

\pnum{5.28}{Reproducible dispersion plot of the delta comb}
With $m=\sh=a=\alpha=1$,
\begin{equation}
E=\frac{k^2}{2},\qquad
\cos q=\cos k+\frac{\sin k}{k}\equiv f(k).
\end{equation}
For each sampled $E\in(0,30]$,
calculate $k=\sqrt{2E}$ and $f(k)$.
If $|f(k)|\le1$, define
$q=\arccos f(k)$ and add the two
states $(q,E)$ and $(-q,E)$.
Otherwise the energy belongs to a
forbidden gap and should not be
plotted as a Bloch eigenvalue.
For the interval $[-\pi,\pi]$ this
constructs the bands in the first
Brillouin zone. The script
\texttt{chapter5/scripts/verify.py}
produces \texttt{figures/bands.pdf}
using $30{,}000$ energy samples.
\checkit{The plotted curve is an implicit
dispersion relation, not the free-electron
parabola; gaps occur where $|f(k)|>1$.}
